
\section{Abstract}
\label{sec-abstract}
Language models have demonstrated remarkable capabilities on standard benchmarks, yet they struggle increasingly from \emph{mode collapse}, the inability to generate diverse and novel outputs.
Our work introduces \textsc{NoveltyBench}, a benchmark specifically designed to evaluate the ability of language models to produce multiple distinct and high-quality outputs.
\textsc{NoveltyBench} utilizes prompts curated to elicit diverse answers and filtered real-world user queries.
Evaluating 20 leading language models, we find that current state-of-the-art systems generate significantly less diversity than human writers.
Notably, larger models within a family often exhibit less diversity than their smaller counterparts, challenging the notion that capability on standard benchmarks translates directly to generative utility.
While prompting strategies like in-context regeneration can elicit diversity, our findings highlight a fundamental lack of distributional diversity in current models, reducing their utility for users seeking varied responses and suggesting the need for new training and evaluation paradigms that prioritize diversity alongside quality.

Code and data: \url{https://novelty-bench.github.io}.
